\addcontentsline{toc}{subsubsection}{Solution 22. The Hidden Quadratic Function}

\begin{dlsummary}[title={\bfseries Solution 22. The Hidden Quadratic Function},breakable]

\begin{proof}
Put $g(x)=\sqrt{f(x)}$. The given inequalities
are equivalent to
\begin{align}
g(x+y)&\ge g(x)+g(y), \tag{1}\\
g(xy)&\le g(x)g(y). \tag{2}
\end{align}
Since $g>0$, condition (1) implies that $g$
is strictly increasing. From (2), $g(1)\ge1$.

\medskip
\noindent\textbf{Step 1. Show that $g(2)=2$.}

By (1),
\[
g(2)\ge2g(1)\ge2,\qquad
g(4)\ge4g(1)\ge4.
\]
The condition $f(f(2))\le16$ gives
\[
g(g(2)^2)\le4.
\]
Since $g(2)^2\ge4$ and $g$ is increasing,
\[
4\ge g(g(2)^2)\ge g(4)
\ge4g(1)\ge4.
\]
Thus $g(1)=1$ and, by strict monotonicity,
\[
\boxed{g(2)=2.}
\]

\medskip
\noindent\textbf{Step 2. Show that $g(x)\le x$.}

For every positive integer $n$, (1) and (2)
give
\[
2^n=2^n g(1)
\le g(2^n)\le g(2)^n=2^n.
\]
Hence
\[
\boxed{g(2^n)=2^n.}
\]

Fix $x>0$ and put
\[
m_n=\left\lfloor\frac{2^n}{x}\right\rfloor.
\]
For sufficiently large $n$, we have $m_n\ge1$.
By monotonicity and (1),
\[
2^n=g(2^n)\ge g(m_nx)\ge m_ng(x).
\]
Consequently,
\[
g(x)\le
\frac{2^n}{\lfloor2^n/x\rfloor}
\xrightarrow[n\to\infty]{}x.
\]
Therefore
\[
\boxed{g(x)\le x\qquad(x>0).}
\]

\medskip
\noindent\textbf{Step 3. Establish equality.}

By (2) and the preceding bound,
\[
1=g(1)
\le g(x)g(1/x)
\le \frac{g(x)}{x}.
\]
Thus $g(x)\ge x$, and hence
\[
g(x)=x\qquad(x>0).
\]
It follows that
\[
\boxed{f(x)=x^2\qquad(x>0).}
\]

Finally, this function satisfies $f(f(2))=16$
and
\[
\sqrt{f(x+y)}
=\sqrt{f(x)}+\sqrt{f(y)}
=\sqrt{f(xy)}
\left(
\frac1{\sqrt{f(x)}}+
\frac1{\sqrt{f(y)}}
\right)
=x+y.
\]
Therefore it is indeed the unique solution.
\end{proof}

\end{dlsummary}